\documentclass[11pt]{article}
\usepackage[margin=1in]{geometry}
\usepackage{amsmath,amssymb,amsthm}
\usepackage{booktabs}
\usepackage{graphicx}
\usepackage{algorithm}
\usepackage{algpseudocode}
\usepackage[numbers,sort&compress]{natbib}
\usepackage[colorlinks=true,allcolors=blue]{hyperref}

\newtheorem{proposition}{Proposition}
\newtheorem{remark}{Remark}

\DeclareMathOperator{\supp}{supp}
\newcommand{\bx}{\mathbf{x}}
\newcommand{\bk}{\mathbf{k}}
\newcommand{\om}{\boldsymbol{\omega}}
\newcommand{\omp}{\om^{\perp}}
\newcommand{\Rmu}{\mathcal{R}_{\mu}}
\newcommand{\bigO}{\mathcal{O}}

\title{A multilevel $\bigO(N\log N)$ algorithm for evaluating the
attenuated Radon transform}
\author{%
  Oren E. Livne\thanks{Pine Birch Analytics, 35 Kelinger Rd, Churchville,
    PA 18966-1033, USA
    (\texttt{oren.livne@gmail.com}, +1 312 533 9130). Corresponding author.}
  \and
  Achi Brandt\thanks{Department of Computer Science and Applied Mathematics,
    The Weizmann Institute of Science, Rehovot 7610001, Israel
    (\texttt{achi.brandt@weizmann.ac.il}).}
}
\date{\today}

\begin{document}
\maketitle

\begin{abstract}
The forward operator of emission tomography is the attenuated Radon transform,
in which the attenuation enters through an exponential whose argument varies
along the line of integration. Evaluating it on an $n^d$ grid for all
significantly different lines costs $\bigO(n^{2d-1})$ operations by direct
quadrature --- one factor of $n$ for each of the $2(d-1)$ parameters of a line,
and one more for traversing it --- and this evaluation is the inner loop of
every iterative reconstruction method. We give an
algorithm that removes the traversal factor, costing
$\bigO(N \log N \log(1/\varepsilon))$ in two dimensions for a prescribed
accuracy $\varepsilon$, and in three dimensions the size of the output itself,
by carrying the recursive length-doubling construction of Brandt
and Dym for plain line integrals over to the attenuated case. Two ingredients
make the transfer work. An exact concatenation identity for a triple of segment
integrals lets the integration length double from level to level with no
approximation. And indexing segments by a starting point on a fixed mesh
together with an integer rise --- so that directions are equispaced in slope and
every segment begins and ends on the mesh --- means that doubling needs no
spatial interpolation at all, only a one-dimensional interpolation in the slope
index, at $\bigO(p)$ operations per integral for an order-$p$ stencil. We
identify the quantities that must be interpolated: the attenuation exponent
rather than its exponential, a segment integral referenced to its own exit point
rather than to its midpoint, and both normalised by arc length. Because the
attenuated merge is nonlinear the interpolation must precede the concatenation,
which is not so for a plain line integral. Constant data is then reproduced
exactly at every level, in two and three dimensions, to $4\times10^{-16}$. The
four-quadrant transform of a $129^2$ image takes $0.63$\,s, and the full
four-dimensional transform of a $33^3$ image --- $2.8\times10^7$ line integrals
--- takes $43$\,s.
\end{abstract}

\section{Introduction}
\label{sec:intro}

\subsection{The forward model}

Emission tomography reconstructs the spatial distribution of a radiotracer
from counts of the photons it emits. Two fields are involved: the activity
$f(\bx)\ge 0$, which is the unknown, and the linear attenuation coefficient
$\mu(\bx)\ge 0$ of the surrounding medium at the photon energy, which is
usually measured separately by a transmission or CT scan.

Parameterise a line in $\mathbb{R}^2$ by a direction and a signed offset,
\begin{equation}
  \om(\theta)=(\cos\theta,\sin\theta),\qquad
  \omp(\theta)=(-\sin\theta,\cos\theta),\qquad
  \bx(s,t)=s\,\omp+t\,\om ,
\end{equation}
and traverse $L(\theta,s)=\{\bx(s,t):t\in\mathbb{R}\}$ in the direction of
increasing $t$. A photon emitted at $\bx(s,t)$ and travelling along $\om$
escapes with probability $\exp(-\int_t^{\infty}\mu)$, so the expected count
along $L(\theta,s)$ is
\begin{equation}
  (\Rmu f)(\theta,s)
  = \int_{\mathbb{R}} f\big(\bx(s,t)\big)\,
    \exp\!\left(-\int_t^{\infty}\mu\big(\bx(s,\tau)\big)\,d\tau\right) dt .
  \label{eq:art}
\end{equation}
This is the attenuated Radon transform, also called the attenuated X-ray
transform. Its distinguishing feature is that the exponent depends on the
emission point $t$, so the attenuation does not factor out of the integral.
Equation~\eqref{eq:art} is the forward model of single-photon emission
computed tomography (SPECT).

Positron emission tomography (PET) detects two collinear annihilation photons
in coincidence. Both must escape, in opposite directions, so the survival
probability is $\exp(-\int_t^{\infty}\mu)\exp(-\int_{-\infty}^{t}\mu)
=\exp(-\int_L\mu)$, independent of $t$, and
\begin{equation}
  (\mathcal{P}_{\mu}f)(\theta,s)
  = \exp\!\left(-\int_L \mu\right)\int_L f .
  \label{eq:pet}
\end{equation}
Classical PET attenuation is thus a multiplicative factor per line. The
algorithm below computes the three line integrals
\begin{equation}
  S=\int_L f,\qquad E=\int_L\mu,\qquad I=(\Rmu f)(\theta,s),
  \label{eq:three}
\end{equation}
and therefore delivers both models: \eqref{eq:art} is $I$, and \eqref{eq:pet}
is $e^{-E}S$. The exponent is genuinely position-dependent only in
\eqref{eq:art}, and that is the case treated here; \eqref{eq:pet} is recovered
as a by-product at no extra cost. Time-of-flight PET, in which the coincidence
is localised along the line and the two one-sided attenuations no longer
combine into a constant, has the structure of \eqref{eq:art} and is covered by
the same machinery.

\subsection{The inverse problem and where evaluation enters}
\label{sec:motivation}

With $\mu$ known, \eqref{eq:art} is a linear Fredholm integral equation of the
first kind for $f$,
\begin{equation}
  \Rmu f = g ,
  \label{eq:inverse}
\end{equation}
where $g$ is the measured sinogram. Novikov \citep{novikov2002} gave an
explicit inversion formula for $\Rmu$, and for the unattenuated case
$\mu\equiv 0$ the classical filtered-backprojection inversion of the Radon
transform applies \citep{natterer2001}. In practice \eqref{eq:inverse} is
solved iteratively rather than by an inversion formula, because the data are
Poisson counts, because the system is large and ill-conditioned, and because
regularisation and physical corrections are easier to impose inside an
iteration.

The standard methods are all of this kind. The expectation-maximisation
algorithm for the Poisson likelihood \citep{sheppvardi1982} updates
\begin{equation}
  f^{(k+1)} = \frac{f^{(k)}}{\Rmu^{*}\mathbf{1}}\;
              \Rmu^{*}\!\left(\frac{g}{\Rmu f^{(k)}}\right),
  \label{eq:mlem}
\end{equation}
a multiplicative --- hence nonlinear --- relaxation; its ordered-subset
variant \citep{hudsonlarkin1994} applies \eqref{eq:mlem} over subsets of
directions in sequence, which makes it a block row-action scheme in the family
of algebraic reconstruction techniques \citep{gordon1970}, i.e.\ a Kaczmarz
relaxation with a nonlinear update. Every sweep of every one of these methods
applies $\Rmu$ once and its adjoint once, so the cost of the reconstruction is
a small multiple of the cost of evaluating $\Rmu$.

The exponent alone understates what is at stake. Direct quadrature costs
$\bigO(n^{2d-1})$ and the algorithm below costs $\bigO(n^{2d-2}\log n)$ in two
dimensions and $\bigO(n^{2d-2})$ in three, so the saving is not the gap between
one power of $N$ and another but a single factor of $n$ --- the number of
samples taken along each line --- reduced to $\bigO(\log n)$ or removed. At
$n = 512$ that factor is $512$, and it multiplies the number of sweeps.

That the projector and its adjoint dominate iterative reconstruction in
practice is not in dispute. On the HRRT brain scanner a CPU reconstruction took
days, which kept a high-sensitivity instrument out of routine use until a GPU
implementation brought one OSEM iteration with sixteen subsets to $16$\,s, a
speed-up of about $125$ \citep{kimye2011}; and projection is enough of the cost
that it is maintained as a separate, dedicated library
\citep{parallelproj2024}. Those are engineering responses --- more arithmetic
per second --- to a cost that is algorithmic. A faster projector composes with
them rather than competing.

Two further considerations motivate a multilevel evaluation specifically
rather than any fast evaluation. First, if $\mu$ is not measured separately
but estimated jointly with $f$ from the emission data, the problem becomes
nonlinear in $\mu$, since $\mu$ sits in the exponent; the natural solver is
then a full approximation scheme (FAS) \citep{brandt1977} in which the
nonlinearity is handled on a hierarchy of grids. Second, a row-action
relaxation of \eqref{eq:inverse} is a smoother in the multigrid sense: it
reduces the components of the residual that the fine grid resolves, and leaves
the smooth components to be handled on coarser grids, so a full multigrid
(FMG) cycle over the image hierarchy is the natural outer solver. The level
hierarchy constructed here --- segment lengths paired with angular resolutions
--- is the same hierarchy such a solver would traverse, which is what makes the
evaluation and the solver composable. We do not develop the solver in this
paper. The present contribution is the forward evaluation; the adjoint, the
nonlinear inversion and the FAS/FMG cycle are the subject of subsequent work
(Section~\ref{sec:discussion}).

\subsection{Related work and contribution}

Fast summation by separating a smooth kernel on a hierarchy of grids is the
subject of the multilevel matrix-multiplication framework of Brandt and
Lubrecht \citep{brandtlubrecht1990} and its extension to oscillatory kernels
\citep{brandt1991}; the fast multipole method \citep{greengardrokhlin1987}
solves the same problem by a different decomposition. For line integrals
specifically, Brandt and Dym \citep{brandtdym1999} gave an $\bigO(N\log N)$
algorithm based on recursive doubling of segment length, and Brandt et al.\
\citep{brandtmann2000} built a multilevel inversion of the unattenuated Radon
transform on it. Fast backprojection by hierarchical decomposition, at
$\bigO(N\log^2 N)$ \citep{basubresler2000}, and transforms built on the
pseudopolar grid \citep{averbuch2008} achieve comparable asymptotics for the
unattenuated transform.

The contribution here is the extension to the attenuated transform
\eqref{eq:art}, where the smoothness of the exponent and of the integrand must
be exploited together. Specifically:
\begin{enumerate}
\item an exact concatenation identity (Proposition~\ref{prop:merge}) for a
  triple of segment integrals, which makes the doubling of segment length
  exact and, in addition, damping: an error committed upstream is multiplied by
  $e^{-E}\le 1$ at every merge;
\item identification of the quantities that admit interpolation --- the
  exponent $E$ rather than $e^{-E}$, and an exit-referenced rather than
  midpoint-referenced attenuated integral --- with the failure mode of each
  alternative;
\item a sampling analysis of the three axes (direction, transverse offset,
  position along the ray) which fixes the level schedule and shows that the
  angular refinement must be placed at the finest levels, not the coarsest;
\item a measured account of the resulting accuracy and cost, including the
  finding that the error budget is dominated not by the angular interpolation
  the method is built around, but by the resampling of the tables in the
  transverse and along-ray variables when the ray frame rotates.
\end{enumerate}
An open-source reference implementation, together with the scripts that
produce every number and figure below, is available
as~\citep{tomogrid}.

\section{Segment integrals and the concatenation identity}
\label{sec:identity}

The algorithm works with integrals over finite segments of a ray rather than
over whole lines. Fix $(\theta,s)$ and write $f(t)$, $\mu(t)$ for the
restrictions to that ray. For $a<b$ define
\begin{equation}
  E[a,b]=\int_a^b\mu(\tau)\,d\tau,\qquad
  S[a,b]=\int_a^b f(t)\,dt,\qquad
  I[a,b]=\int_a^b f(t)\,e^{-\int_t^b\mu(\tau)d\tau}\,dt .
  \label{eq:triple}
\end{equation}
The attenuation in $I$ is referenced to the segment's own exit point $b$, so
every weight lies in $(0,1]$. This choice is what makes the recursion stable
(Remark~\ref{rem:reference}).

\begin{proposition}[Concatenation]
\label{prop:merge}
For any $c\in(a,b)$,
\begin{equation}
  E[a,b]=E[a,c]+E[c,b],\qquad
  S[a,b]=S[a,c]+S[c,b],\qquad
  I[a,b]=e^{-E[c,b]}\,I[a,c]+I[c,b].
  \label{eq:merge}
\end{equation}
\end{proposition}

\begin{proof}
The first two are additivity of the integral. For the third, split at $c$. For
$t\in[a,c]$ write $\int_t^b\mu=\int_t^c\mu+\int_c^b\mu$; the second term does
not depend on $t$ and factors out as $e^{-E[c,b]}$, leaving $I[a,c]$. For
$t\in[c,b]$ the integrand is unchanged and contributes $I[c,b]$.
\end{proof}

Equivalently, the lower-triangular matrices
\begin{equation}
  M[a,b]=\begin{pmatrix} e^{-E[a,b]} & 0\\ I[a,b] & 1\end{pmatrix}
  \qquad\text{satisfy}\qquad
  M[a,b]=M[a,c]\,M[c,b],
\end{equation}
so the triple $(S,E,I)$ carries a semigroup structure under concatenation.
Identity \eqref{eq:merge} involves no quadrature and no interpolation: given
the triples of two adjacent segments, the triple of their union is obtained
exactly. This is what permits the segment length to double from level to level
at no cost in accuracy.

\begin{remark}[Why $I$ is referenced to the exit point]
\label{rem:reference}
Referencing the attenuation to the segment midpoint $m$, i.e.\ using
$\tilde I[a,b]=\int_a^b f(t)e^{-\int_t^m\mu}dt$, gives weights as large as
$e^{+E[a,m]}$, which grow without bound with segment length and attenuation
strength. With the exit-referenced definition \eqref{eq:triple} all weights are
in $(0,1]$, and \eqref{eq:merge} shows that an error $\varepsilon$ in $I[a,c]$
propagates into $I[a,b]$ as $e^{-E[c,b]}\varepsilon$. Strong attenuation
therefore improves the conditioning of the recursion. This is confirmed in
Section~\ref{sec:exact}, where strong attenuation never degrades the computed
transform.
\end{remark}

\begin{remark}[Why $E$ and not $e^{-E}$ is interpolated]
\label{rem:exponent}
$E$ is an additive line integral of $\mu$, smoothed along the ray by the
averaging in \eqref{eq:triple}, with derivatives bounded in terms of those of
$\mu$. Its exponential spans many orders of magnitude --- in the experiments
below $E$ reaches $12$, so $e^{-E}\approx 6\times 10^{-6}$ --- and has
correspondingly large relative derivatives. The algorithm interpolates $E$ and
applies the exponential afterwards. This is the sense in which smoothness is
used twice: once in the exponent, through $E$, and once in the integrand,
through $I$.
\end{remark}

\section{Resolution requirements}
\label{sec:sampling}

Three rules fix how densely the integrals must be computed; all three are in
\citep{brandtdym1999} and all three are needed to see why the number of
integrals per level does not depend on the length.

Write $F(\bx, L, \theta)$ for the integral over a segment of length $L$,
normalised by $L$, so that $F$ is the \emph{mean} of the data along the segment.
Perturbing the direction,
\begin{equation}
  F(\bx, L, \theta + \xi) = F(\bx, L, \theta) + R,
  \qquad |R| \le \tfrac{\xi L}{4}\max|\nabla g| ,
  \label{eq:angcrit}
\end{equation}
so the angular resolution needed for a given accuracy is proportional to $L$:
\emph{double the length, double the number of directions}. Moving the segment a
distance $\epsilon$ along its own direction changes $F$ by at most
$2\epsilon\max|g|/L$, so the spacing of evaluation points along the direction is
inversely proportional to $L$: \emph{double the length, halve the number of
positions}. Transverse to the direction no averaging occurs, so the spacing
there is set by the bandwidth of the data and is the same at every level. The
three together give the property the method rests on:
\begin{equation}
  \text{number of integrals at length } L \;=\; \text{independent of } L .
  \label{eq:flat}
\end{equation}
Finally $\partial F/\partial L$ is $\bigO(1/L)$, so $L$ itself need only be
sampled geometrically --- one level per doubling.

The same counting in $d$ dimensions gives directions growing as $(L/h)^{d-1}$
and positions falling as $h/L$, so in three dimensions the table grows by two
per level and the total is dominated by the top level. That is not a defect:
the three-dimensional X-ray transform is a four-dimensional object, and the top
level \emph{is} the output.

\section{The algorithm}
\label{sec:algorithm}

\subsection{Indexing}

Everything turns on how segments are indexed. Let $H$ be the data mesh and
$h \le H$ the integration mesh transverse to the sweep; their ratio
\begin{equation}
  \sigma = H/h \ \ (\ge 1)
  \label{eq:sigma}
\end{equation}
is the single parameter controlling how finely the tables are sampled relative
to the data, and is the paper's $H/h$. It is the main accuracy knob: the error
of one interpolation is $\bigO(\sigma^{-p})$ and the work is
$\bigO(\sigma^{2(d-1)})$. For the family of
directions within $45^\circ$ of $+x$, a segment of horizontal extent $L$ runs
from a mesh point $(x_i, y_j)$ to $(x_i + L,\, y_j + m h)$ with $m$ an
\emph{integer}. Its slope is $mh/L$, so directions are equispaced in $\tan\theta$
and the slope set at level $k$ has spacing $h/L_k$, doubling in size exactly as
\eqref{eq:angcrit} requires. Crucially, every segment both begins and ends on
the mesh. Level $k$ carries
\begin{equation}
  L_k = 2^k L_0, \qquad
  x\text{ starts spaced } L_k, \qquad
  y\text{ starts spaced } h, \qquad
  m \in [-R_k, R_k],\; R_k = L_k/h,
\end{equation}
whose product is independent of $k$, which is \eqref{eq:flat}.

\subsection{Doubling}

For a level-$(k+1)$ target of extent $2L$ and rise $M$:

\paragraph{$M$ even.} The direction already exists at level $k$. The two halves
are the level-$k$ segments $[2i,\, j,\, M/2]$ and $[2i+1,\, j + M/2,\, M/2]$ ---
the second beginning exactly where the first ends, both on the mesh. No
approximation whatever.

\paragraph{$M$ odd.} The direction is new. Interpolate the \emph{slope index}
over $p$ neighbouring integer rises, separately for each half, at the same
starting point:
\begin{equation}
  A = \sum_a w_a\, T_k[2i,\; j,\; m_a], \qquad
  B = \sum_a w_a\, T_k[2i+1,\; j + m_a,\; m_a],
  \label{eq:double}
\end{equation}
with $w_a$ the midpoint Lagrange weights for the half-rise $M/2$. The start row
of the second half moves \emph{with} the slope index $m_a$, because each
direction's second half begins where its own first half ended. That is what
removes any need for spatial interpolation: \eqref{eq:double} contains only a
one-dimensional interpolation in $m$ and an integer shear in $(j,m)$, at
$\bigO(p)$ operations per integral.

\subsection{What is interpolated}

Three choices, each of which the measurements in Section~\ref{sec:results}
penalise if made differently.

\emph{Normalise by arc length.} A segment of extent $L$ and rise $mh$ has arc
length $\sqrt{L^2 + (mh)^2}$, so interpolating an unnormalised integral in the
slope index interpolates the arc length along with it and leaves a purely
geometric error unrelated to the data. Storing mean values removes it: a
constant field is then reproduced exactly.

\emph{Interpolate $E$, not $e^{-E}$}, and exponentiate afterwards, for the
reason given in Remark~\ref{rem:exponent}; and \emph{reference $I$ to the
segment's own exit point}, for the reason given in
Remark~\ref{rem:reference}.

\emph{Interpolate, then concatenate.} For a plain line integral the merge is an
average, so it commutes with the interpolation and the order is immaterial ---
both readings of \citep{brandtdym1999} give identical results. The attenuated
merge is nonlinear in $E$ and does not commute. Each half's triple must be
interpolated first and the two concatenated afterwards by \eqref{eq:merge},
which keeps the concatenation exact. Writing $\ell$ for the arc length of each
half of the \emph{target} segment,
\begin{equation}
  \bar S = \tfrac12(\bar S_A + \bar S_B), \quad
  \bar E = \tfrac12(\bar E_A + \bar E_B), \quad
  \bar I = \tfrac12\big(e^{-\bar E_B \ell}\, \bar I_A + \bar I_B\big).
  \label{eq:merge_norm}
\end{equation}

\subsection{Base level, families, and three dimensions}

Only the base level touches the data: the shortest segments are evaluated by
Gauss--Legendre quadrature with the fields sampled by tensor-product
interpolation from the image grid. Four sweeps are needed in two dimensions and
six in three --- the shallow cone about each axis direction --- because the
attenuated transform is \emph{directed}: the exponent is referenced to the exit
end, so $\theta$ and $\theta + \pi$ differ. Each is the same code applied to a
reflected or transposed image.

In three dimensions a segment carries two integer rises, and the construction is
unchanged. The targets form a grid in the two slope indices, so the
two-dimensional interpolation separates into two one-dimensional passes and
still costs $\bigO(p)$ per entry rather than $\bigO(p^2)$; the shear becomes two
successive integer gathers, one per transverse axis.

\subsection{Complexity}
\label{sec:complexity}

With \eqref{eq:flat} the table size is the same at every level, there are
$l \approx \log n$ levels, and \eqref{eq:double} costs $\bigO(p)$ per integral,
so in two dimensions
\begin{equation}
  W = \bigO(p\,\sigma^2 N \log N) .
\end{equation}
The error of one interpolation is $\bigO((h/H)^p) = \bigO(\sigma^{-p})$; over
$l$ levels, reaching a prescribed accuracy $\varepsilon$ requires
$l\,\sigma^{-p} \approx \varepsilon$, so at fixed $\sigma$
\begin{equation}
  p = \bigO\!\big(\log(1/\varepsilon) + \log\log n\big),
  \qquad\text{hence}\qquad
  W = \bigO\!\big(N \log N \log(1/\varepsilon)\big).
  \label{eq:cost}
\end{equation}
Optimising $p$ against $\sigma$ instead gives the form
$\bigO(p\,n(\log n)^{1+2/p})$ reported in \citep{brandtdym1999}, minimised at
$p = 2\ln\log n$ to give $\bigO(n \log n \log\log n)$. Either way the cost per
integral is a constant times the interpolation order, which is what
distinguishes this indexing from one in which the ray frame rotates with the
direction: there each entry would need a scattered $d$-dimensional resample at
$\bigO(p^d)$, and in three dimensions that is more expensive than computing the
entry from the image directly.

\section{Numerical results}
\label{sec:results}

Two codes are exercised. \texttt{bdlines} reproduces
\citep{brandtdym1999} for plain line integrals, so that the discretisation, the
merging, the interpolation and the work count can each be checked against
published numbers before any of it is carried into the attenuated case;
\texttt{tomogrid} is the attenuated transform on the same indexing. All numbers
below are produced by committed scripts over committed data.

\subsection{Exactness where the method claims it}
\label{sec:exact}

The properties that must hold to roundoff are the sharpest test of the
indexing, and each of them failed for a formulation in which the ray frame
rotates with the direction.

A direction already present at the base level is merged and never interpolated,
and the normalised trapezoid rule composes exactly, so the recursion must agree
with the direct discretisation to roundoff: measured $1.2\times10^{-16}$.
Constant data must be reproduced exactly at every level, which additionally
tests the arc-length normalisation: measured $4.4\times10^{-16}$, in two and in
three dimensions. With $\mu\equiv0$ the computed $I$ equals $S$ to
$10^{-12}$; with $f\equiv0$ both vanish identically while $E$ does not; and the
transform is linear in $f$ to $10^{-12}$.

\subsection{Reproducing Brandt and Dym}
\label{sec:reproduce}

The central claim of \citep{brandtdym1999} is that halving the integration mesh
reduces the approximation error by $2^p$. Measured over $H/h = 1,2,4$ on the
paper's sinusoidal target: linear doubling gives $3.8\times$ and $3.9\times$,
cubic gives $16\times$ and $19\times$, quintic gives $2^6$ until it reaches the
discretisation floor. The discretisation errors agree with the published values
to within a few per cent --- $0.085\%$ against $0.092\%$, $0.019\%$ against
$0.025\%$, $0.0062\%$ against $0.008\%$ --- as does the near-independence of the
approximation error from $L$.

The absolute approximation error with \emph{linear} doubling comes out about
$3.5\times$ larger than published, with the same rates. Almost all of the
difference is the interpolation of the second half of each new direction, whose
start row moves with the slope index; the first half contributes about five
times less. Both readings of the text were implemented and give identical
results, because with linear interpolation the interpolation and the averaging
commute, so the published Figure 3(b) is needed to settle which four length-$L$
integrals are averaged. The discrepancy does not affect the design criterion,
which is that the approximation need only be as accurate as the discretisation
it approximates: with cubic doubling the error at $H/h=1$ is $0.015\%$, already
well below the $0.085\%$ discretisation error.

\subsection{The attenuated transform in two dimensions}

Table~\ref{tab:spect2d} reports the attenuated transform of a $129^2$ image over
all four direction families, against a dense-trapezoid evaluation of the
analytic phantom, line by line on the segment endpoints the algorithm produces.

\begin{table}[t]
\centering\small
\caption{Attenuated transform, $129^2$ image, all four direction families,
against the continuum. $\sigma = H/h$; $p$ is the slope-interpolation order.
Operation counts are for the whole transform.}
\label{tab:spect2d}
\begin{tabular}{rr rrr}
\toprule
$p$ & $\sigma$ & $\|I - I_{\mathrm{ref}}\|_\infty / \max|S|$ & Mflop & seconds\\
\midrule
$2$ & $1$ & $3.06\!\times\!10^{-2}$ & $38.2$ & $0.59$\\
$2$ & $2$ & $9.73\!\times\!10^{-3}$ & $142.9$ & $1.57$\\
$2$ & $4$ & $3.27\!\times\!10^{-3}$ & $552.1$ & $9.00$\\
\addlinespace
$4$ & $1$ & $1.04\!\times\!10^{-2}$ & $47.7$ & $0.38$\\
$4$ & $2$ & $2.26\!\times\!10^{-3}$ & $180.8$ & $1.28$\\
$4$ & $4$ & $1.06\!\times\!10^{-3}$ & $703.4$ & $9.51$\\
\addlinespace
$6$ & $1$ & $2.81\!\times\!10^{-2}$ & $57.2$ & $0.56$\\
$6$ & $2$ & $2.25\!\times\!10^{-3}$ & $218.7$ & $1.63$\\
$6$ & $4$ & $1.01\!\times\!10^{-3}$ & $854.7$ & $9.58$\\
\bottomrule
\end{tabular}
\end{table}

Raising the order is cheaper than refining the mesh, as \eqref{eq:cost} says it
should be: at $\sigma = 1$, going from $p=2$ to $p=4$ reduces the error by a
factor of three for a quarter more work, whereas reaching the same error by
refining $\sigma$ costs four times the work.  At $\sigma = 1$ the sixth-order
stencil is \emph{worse} than the fourth ($2.81\times10^{-2}$ against
$1.04\times10^{-2}$), for the reason returned to in Section~\ref{sec:3d}; by
$\sigma = 4$ the two agree to within $5\%$, both having reached the
discretisation floor.

Table~\ref{tab:refine2d} refines the grid at fixed parameters, which is where
\eqref{eq:cost} is tested directly.

\begin{table}[t]
\centering\small
\caption{Refinement at $p=4$, $\sigma=2$, all four families. ``lines'' is the
number of attenuated line integrals produced. Each row doubles $n$, so $N$
quadruples: a pure $\bigO(N)$ method would multiply the work by exactly four and
an $\bigO(N^{3/2})$ one by eight. The error at $n=513$ is measured on every
fourth offset and every eighth slope; the reference costs far more than the
transform it checks.}
\label{tab:refine2d}
\begin{tabular}{rrr rrr}
\toprule
$n$ & error & order & lines & Mflop & work ratio\\
\midrule
$33$  & $3.67\!\times\!10^{-2}$ & ---    & $33\,540$    & $9.2$   & ---\\
$65$  & $1.21\!\times\!10^{-2}$ & $1.60$ & $132\,612$   & $41.0$  & $4.46$\\
$129$ & $2.26\!\times\!10^{-3}$ & $2.42$ & $527\,364$   & $180.8$ & $4.41$\\
$257$ & $2.95\!\times\!10^{-4}$ & $2.94$ & $2\,103\,300$ & $791.1$ & $4.37$\\
$513$ & $1.70\!\times\!10^{-5}$ & $4.12$ & $8\,400\,900$ & $3438.4$ & $4.35$\\
\bottomrule
\end{tabular}
\end{table}

The work ratios are $4.46$, $4.41$, $4.37$, $4.35$, against the $4$ of a linear
method and the $8$ of direct quadrature: this is $\bigO(N \log N)$ measured, the
logarithm visible as the drift above four, itself decreasing because the level
count grows only as $\log n$. The error converges at an order rising through
$1.60$, $2.42$, $2.94$ to $4.12$ --- the fourth order of the cubic image
interpolant, reached. At $n = 513$ the transform delivers $8.4$ million
attenuated line integrals in $54$\,s.

Against direct quadrature on the same lines, with the same number of samples
per line as the grid has cells, the measured ratios are $2.7$, $6.0$ and $7.9$
at $n = 65$, $129$ and $257$ --- roughly doubling with $n$, which is the
$n/\log n$ of Section~\ref{sec:complexity} rather than any difference of
exponents.

\subsection{Three dimensions}
\label{sec:3d}

The full four-dimensional transform of a $33^3$ image --- six direction
families, $2.8\times10^7$ line integrals --- takes $43$\,s. Separating the two
error sources on one family, against a reference built from the same image
interpolant and against the continuum:

\begin{center}\small
\begin{tabular}{rr rrr}
\toprule
$p$ & $\sigma$ & $e_{\mathrm{alg}}$ & $e_{\mathrm{disc}}$ & seconds\\
\midrule
$4$ & $1$ & $7.08\!\times\!10^{-2}$ & $3.55\!\times\!10^{-2}$ & $6.4$\\
$4$ & $2$ & $7.16\!\times\!10^{-3}$ & $2.85\!\times\!10^{-2}$ & $143.6$\\
$6$ & $1$ & $1.04\!\times\!10^{-1}$ & $3.55\!\times\!10^{-2}$ & $6.6$\\
\bottomrule
\end{tabular}
\end{center}

Two dimensions meets the design criterion at $\sigma = 1$; three dimensions
needs $\sigma = 2$, where the algorithmic error is four times \emph{below} the
discretisation error, having fallen tenfold from $\sigma=1$. Raising the order
instead does not work at $\sigma = 1$: $p=6$ is worse than $p=4$. On equispaced
nodes the Lebesgue constant of Lagrange interpolation grows like $2^p$, so the
order only pays once the table spacing is comfortably inside the field's Nyquist
rate, which at $\sigma = 1$ it is not. This is the practical content of
\eqref{eq:cost}: $\varepsilon$ is bought by the order, but only above a
threshold in $\sigma$ that depends on the dimension.

\subsection{Against a rotating ray frame}

The same attenuated transform was first implemented with tables indexed by
$(\theta, s, t)$ in a ray frame that rotates with the direction. Refining the
direction then moves every sample point, so each entry requires a scattered
$d$-dimensional resample: $192$ gathers per entry in two dimensions and $3072$
in three, against $2\,g\,p_{\mathrm{img}}^d = 256$ and $768$ to compute the
entry from the image directly. In two dimensions that is marginally profitable;
in three it is four times more expensive than not using the hierarchy at all,
and the crossover against direct evaluation sits near $n\approx200$ rather than
$n = 65$. Measured, the $129^2$ transform took $15.4$\,s for a quarter of the
directions, against $0.63$\,s here for all of them. The lesson is that the cost
and the error budget of this method are both decided by the indexing, not by the
interpolation machinery laid over it.

\section{Discussion}
\label{sec:discussion}

\paragraph{Limitations.}
Four, stated precisely. (i) Directions come out equispaced in slope rather than
in angle, and the transverse spacing between neighbouring lines therefore varies
with angle, by $\sqrt2$ across a quadrant; \citep{brandtdym1999} notes both an
equispaced-grid variant and the simpler remedy of interpolating to a uniform
grid at the final stage only. (ii) On rays of very large total attenuation the
true signal underflows, and what the method controls there is a small error
\emph{absolute} to $\max|S|$, not a small relative error; carrying the mean
survival weight $I/S$ instead of $I$ would address this at the cost of an
indeterminate form wherever a segment misses the support. (iii) Zero extension
outside the tables is exact only for compactly supported fields. (iv) The order
$p$ buys accuracy only above a threshold in $\sigma$ that grows with dimension,
because equispaced Lagrange stencils have Lebesgue constants growing like $2^p$;
better-conditioned transfers would remove the threshold and are the natural way
to make \eqref{eq:cost} hold uniformly.

\paragraph{The adjoint, and the inverse problem.}
The adjoint $\Rmu^{*}$, which every method in Section~\ref{sec:motivation}
applies once per sweep alongside the forward operator, is the transpose of the
construction step by step: the concatenation transposes to a scatter and the
slope interpolation to anterpolation, in the sense of
\citep{brandtlubrecht1990}. Because the doubling here is an indexed
multiply--add rather than a scattered resample, its transpose is again an
indexed multiply--add, and the adjoint costs what the forward operator costs.
That is what makes the present formulation, rather than the rotating-frame one,
a usable basis for the reconstruction itself.

Beyond that lies the inverse problem sketched in
Section~\ref{sec:motivation}: the linear solve for $f$ with $\mu$ known, and
then the joint estimation of $(f,\mu)$, in which $\mu$ enters nonlinearly
because it sits in the exponent. A row-action relaxation of \eqref{eq:inverse}
smooths the error components the fine grid resolves and leaves the rest to
coarser grids, so a full multigrid cycle over the image hierarchy is the natural
outer solver, and a full approximation scheme \citep{brandt1977} the natural
treatment of the nonlinearity. The level structure built here --- lengths paired
with angular resolutions --- is the hierarchy such a solver would traverse,
which is what makes evaluation and solver composable. \citep{brandtdym1999}
notes that its own construction can in principle be inverted directly; whether
that carries over to the attenuated case is open.

\section{Conclusion}

The attenuated Radon transform admits an $\bigO(N\log N)$ multilevel
evaluation. The construction rests on an exact concatenation identity for the
triple $(S,E,I)$, on interpolating the exponent rather than its exponential and
the exit-referenced rather than the midpoint-referenced attenuated integral,
and on a level schedule that refines the angular grid at the finest levels
where segments are short. On a battery of attenuation fields the error the
schedule introduces stays an order of magnitude below the discretisation error
of the grid and converges with it, at a cost already several times below direct
evaluation at moderate resolution and improving as $\bigO(\sqrt{N}/\log N)$.

\subsection*{Code}

The reference implementation, the reproduction of \citep{brandtdym1999}, the
test suite and the scripts producing every number above are
at~\citep{tomogrid}; its README lists the commands.

\bibliographystyle{plainnat}
\bibliography{refs}

\end{document}
